% Pursuing Stacks (cotes 134-2 to 134-8): archive cote 134-8, site page 30 (PDF page 31), typescript page 584.
% typescript. Draft compared: ps7.tex:1392-1455.
% AI-transcribed by Claude Opus 5.5 (claude-opus-5-5, adaptive thinking, effort high), Message Batches API, batch ps4 (msgbatch_012Aoc1NuskATPsyBeKSYK7J), 2026-11-12; reviewed against both supplied scan halves; see CHECK_DECISIONS.tsv.
\scanpage{134-8/30}
\mspage{584}

where now ${}_{b\hand{\backslash}}A$  is the category of pairs
\[ (a,u) \ \text{ with } \ u: b \hand{\to} f_{o}(a) \quad . \]
\over{\unsure{We}}{We}'ll have to assume npw that both sheaves of rings $\underline{O}_A$ , $\underline{O}_B$  are con-%
stant,  and correspond to the same ring k.  \xout{The commutation of f} Thus
Modules on A and B are just contravariant functors from these catego-%
ries to the category  $Ab_k$  of \xout{all} k-modules,  and accordingly $f_!$
(\xout{dual} \inserted[typed]{left adjoint \xout{to}} to the composition functor $f^*$ )  may be computed by a well known
formula,  involving direct limits on the categories ${}_{b\hand{\backslash}}A$ :
\[ (7) \qquad f_!(F)(b) \hand{\simeq} \underset{b\hand{\backslash} A}{\underset{\hand{\longrightarrow}}{\lim}} F(a) \qquad , \]
where the limit in the second member is relative to the \xout{funct} composi-%
tion
\[ {}_{b\hand{\backslash}}A \hand{\longrightarrow} A \overset{\ F\ }{\hand{\longrightarrow}} Ab_k \quad . \]
(NB This formula is dual to the formula for $f_*$ ,  the right adjoint of
$f^*$ ,  closer to intuition  - to mine at any rate -  because the $\underset{\hand{\longleftarrow}}{\lim}$ in
\[ f_*(F)(b) \hand{\simeq} \underset{A_{\hand{/}b}}{\underset{\hand{\longleftarrow}}{\lim}} F(a) \]
may be ``visualized'' as \deleted{\xout{\illegible{1}sections\illegible{1}}}  the \xout{set of} set of ``sections'' of
F over $A_{/b}$ \over{\unsure{,}}{.}) From this formula (7) follows at once commutation of
$f_!$  with colocalization\xout{t}.  To get the corresponding result for $Lf_!$ ,
we have only to use the fact that the relevant inverse image functor
(corresponding to  ${}_{b\hand{\backslash}}A \hand{\to} A$ ) transforms projective\deleted{\xout{\illegible{1}}} modules into
% NOTE: from here to the foot of the page the text is typed on a strip pasted over the page; the older typed text underneath shows through faintly and is not transcribed
projective modules.  As the map  ${}_{b\hand{\backslash}}B \hand{\longrightarrow} B$  is a \uline{cofibering functor}
\uline{with discrete fibers},  so is the map
\[ h: A' = {}_{b\hand{\backslash}}A \hand{\longrightarrow} A \]
deduced by base change.  Now for any \inserted{such} functor  $h: A' \hand{\to} A$ between small
categories,  the inverse image functor
\[ h^* : A^{\overset{\hand{\illegible{1}}}{\wedge}} \hand{\longrightarrow} A'^{\wedge} \ , \quad \text{or} \quad h^*_k : A^{\wedge}_k \hand{\to} A'^{\wedge}_k \]
carries \inserted[typed]{indeed} projectives into projectives.  This statement is dual formally
to the corresponding statement for injectives,  valid when we make on
h the dual assumption of being \uline{fibering with discrete fibers}  -  in the
latter case the \hand{(}well known\hand{)} proof comes out formally from the fact
that the left adjoint functor $h_!$ or $h^k_!$ carries monomorphisms into
monomorphisms \deleted{i.e. is exact}  -  a fact which we used in section 135
\margin{\tikz{\draw[thin] (0,0)--(1.1,0);}}in the form  $Lh^k_! = h^k_!$ (in case $k=\underline{Z}$).  In the present case the proof is
essentially the dual one  -  \over{as}{as} a matter of fact, as there are enough

\margin{30}

% Visual collation: both supplied image halves inspected; see CHECK_DECISIONS_134_7_8.tsv.
